%% hanging-punct.tex — clreq 行尾点号悬挂（直排）
%% 【测试配置】ltc-cn-vbook + \标点设置{punct-hanging=true}
%% clreq 注明悬挂适用于直排（横排港台式点号居中不宜）。
%% 数字串长度递增扫过列末各相位，其中若干段的点号恰好落在列末：
%% 开启悬挂后该点号整幅移出列内，字形挂在版口之外，列内正文不被压缩。
%% 关闭 punct-hanging 时同一列的点号要挤进列内、整列字距被压缩，
%% 是断言「悬挂列步长未被压缩」的负对照（见 clreq_test.py）。
\documentclass{ltc-cn-vbook}
\setmainfont{TW-Kai}

\标点设置{punct-hanging=true}

\begin{document}

\begin{正文}
一二三四五六七八九十一二三四五六七八九十一二三四五六七八九十一二，悬甲

一二三四五六七八九十一二三四五六七八九十一二三四五六七八九十一二三，悬乙

一二三四五六七八九十一二三四五六七八九十一二三四五六七八九十一二三四，悬丙

一二三四五六七八九十一二三四五六七八九十一二三四五六七八九十一二三四五。悬丁

一二三四五六七八九十一二三四五六七八九十一二三四五六七八九十一二三四五六。悬戊

一二三四五六七八九十一二三四五六七八九十一二三四五六七八九十一二三四五六七。悬己

一二三四五六七八九十一二三四五六七八九十一二三四五六七八九十一二三四五六七八，悬庚

一二三四五六七八九十一二三四五六七八九十一二三四五六七八九十一二三四五六七八九，悬辛
\end{正文}

\end{document}
